\section{Key Rotation}

\label{sec:rotation}
%% =====================================================================

\begin{definition}[Proactive Resharing]
Given an existing $t$-of-$n$ sharing of key $x$ on polynomial $f$,
a reshare protocol produces a new $t'$-of-$n'$ sharing on a fresh
polynomial $g$ with $g(0) = f(0) = x$.  The new shares are computed
by having each old signer contribute an additive share of zero
(a random polynomial $h_i$ with $h_i(0) = 0$).  The new shares are
$g(j) = \sum_{i \in S} \lambda_i f(i) + \sum_{i \in S} h_i(j)$ for
each new party $j$.
\end{definition}

\begin{theorem}[Reshare Safety]
\label{thm:reshare}
Proactive resharing satisfies:
\begin{enumerate}[nosep]
  \item \textbf{Public key preservation.}  The aggregate public key
        $X = x \cdot G$ is unchanged after reshare.
  \item \textbf{Old share invalidation.}  Old shares lie on polynomial
        $f$; new shares lie on polynomial $g \neq f$.  Old shares
        cannot be combined with the new protocol state.
  \item \textbf{No key reconstruction.}  The secret $x$ is never
        materialized during the reshare.  Each party sees only its
        own old share plus received shares-of-zero.
\end{enumerate}
\end{theorem}

\begin{proof}
(1) follows because $g(0) = f(0) = x$, so $x \cdot G$ is unchanged.
(2) follows because $g$ is a uniformly random polynomial conditioned on
$g(0) = x$, so the evaluations $g(j)$ are independent of $f(i)$ for
$j \neq i$.  Mixing old and new shares produces interpolation over
inconsistent polynomials, which yields an incorrect secret.
(3) follows because each old signer generates their share-of-zero
locally (using fresh randomness), and the zero-secret polynomial
$h_i$ has $h_i(0) = 0$.  The sum $\sum h_i(j)$ re-randomizes the
polynomial without any party knowing the full randomness.
\end{proof}

Key rotation is performed periodically (default: every 7 days) to
limit the exposure window.  An attacker must compromise $t$ shares
within a single rotation epoch; shares from different epochs are
incompatible.

%% =====================================================================
